\begin{myexercise}
    \section*{📐 Aufgabenstellung: Hierarchie der 2D-Formen}

    In dieser Aufgabe erstellen Sie eine Klassenhierarchie für zweidimensionale Formen in Python.

    \subsection*{1. Die Basisklasse \texttt{ZweiDForm}}

    \begin{enumerate}
        \item \textbf{Erstellen Sie die Basisklasse \texttt{ZweiDForm}.}
        \begin{itemize}
            \item \textbf{Klassenattribut:} Definieren Sie das Klassenattribut \texttt{dimension} und setzen Sie es auf den Wert \texttt{2}.
            \item Implementieren Sie den \textbf{Konstruktor (\texttt{\_\_init\_\_(self, name)})}. 
            \item \textbf{Methoden-Stubs:} Implementieren Sie \texttt{berechne\_flaeche(self)} und \\\texttt{berechne\_umfang(self)}. Diese Methoden sollen jeweils eine \textbf{Warnmeldung} ausgeben (mittels \texttt{print()}), dass sie in den Unterklassen implementiert werden müssen.
        \end{itemize}

        \item \textbf{Testen Sie die Basisklasse.}
        \begin{itemize}
            \item Führen Sie den folgenden Code aus und stellen Sie sicher, dass der Output exakt dem erwarteten Output entspricht:
            \begin{lstlisting}[language=Python, basicstyle=\ttfamily\small]
                quadrat = ZweiDForm("Quadrat")
                print(f"Ein {quadrat.name} hat {quadrat.dimension} Dimensionen.")
                quadrat.berechne_flaeche()
                quadrat.berechne_umfang()
            \end{lstlisting}
            \textbf{Erwarteter Output:}
            \begin{lstlisting}
                Ein Quadrat hat 2 Dimensionen.
                ACHTUNG: Die Methode 'berechne_flaeche' muss in der Unterklasse Quadrat definiert werden.
                ACHTUNG: Die Methode 'berechne_umfang' muss in der Unterklasse Quadrat definiert werden. 
            \end{lstlisting}
        \end{itemize}
    \end{enumerate}

    \subsection*{2. Konkrete 2D-Formen}

    \begin{enumerate}
        \setcounter{enumi}{2} % Zähler fortsetzen
        \item Erstellen Sie zwei Unterklassen von \texttt{ZweiDForm}: \textbf{\texttt{Kreis}} und \textbf{\texttt{Rechteck}}.
            \begin{itemize}
                \item \textbf{Konstruktorbestimmung:} \textbf{Erstellen} Sie für jede Klasse den \textbf{Konstruktor} und \textbf{bestimmen} Sie dabei, welche \textbf{Instanzattribute} benötigt werden, um den Umfang und die Fläche zu berechnen.
            \end{itemize}

        \item \textbf{Implementierung der Methoden.} Überschreiben Sie die Methoden:
            \begin{itemize}
                \item \textbf{Für \texttt{Kreis}:}
                    \begin{itemize}
                        \item \texttt{berechne\_flaeche()}: $A = \pi \cdot r^2$.
                        \item \texttt{berechne\_umfang()}: $U = 2 \cdot \pi \cdot r$.
                    \end{itemize}
                \item \textbf{Für \texttt{Rechteck}:}
                    \begin{itemize}
                        \item \texttt{berechne\_flaeche()}: $A = a \cdot b$.
                        \item \texttt{berechne\_umfang()}: $U = 2 \cdot (a + b)$.
                    \end{itemize}
            \end{itemize}

        \item \textbf{Testen Sie die Implementation.}
            \begin{itemize}
                \item Testen Sie dass für einen Kreis mir Radium $3\,\text{cm}$ der Umfang und die Fläche korrekt berechnet werden.
                \item Testen Sie, dass für ein Rechteck mit Seitenlängen $4\,\text{cm}$ und $6\,\text{cm}$ die Fläche und der Umfang korrekt berechnet werden.
            \end{itemize}
    \end{enumerate}

    \subsection*{3. Das allgemeine N-Eck}

    \begin{enumerate}
        \setcounter{enumi}{5} % Zähler fortsetzen
        \item Fügen Sie die Unterklasse \textbf{\texttt{RegelmaessigesNEck}} hinzu (von \texttt{ZweiDForm} abgeleitet).
        \begin{itemize}
            \item \textbf{Konstruktorbestimmung:} \textbf{Erstellen} Sie den \textbf{Konstruktor}, welcher die Seitenlänge und die Anzahl der Ecken als Attribut speichert. 
            \item \textbf{Implementieren} Sie \texttt{berechne\_flaeche()}: $$A=\frac{n\cdot a^2}{4\tan{\frac{180^\circ}{n}}}$$
            \item \textbf{Implementieren} Sie \texttt{berechne\_umfang()}: $U = n \cdot a$.
        \end{itemize}

        \item \textbf{Testen Sie die Klasse \texttt{RegelmaessigesNEck}}.
    \end{enumerate}

\end{myexercise}